\chapter{The Public-Key Idea and Diffie-Hellman}

In Chapter 17 we saw how to harvest physical entropy and stretch it into unpredictable random keys. Yet throughout Parts I and II, every system we constructed rested on an unstated assumption: Alice and Bob already share a secret. The one-time pad needs a shared pad. AES needs a shared key. HMAC needs a shared MAC key. Authenticated encryption needs both. Every one of these schemes is a machine that turns a small shared secret into confidentiality and integrity for a large conversation. None of them tells you how the shared secret got there in the first place.

This gap was the central practical problem of twentieth-century cryptography, and for most of that century the only known answer was depressing: meet in person, or trust a courier. If you want to talk securely with $n$ people, you need $n$ separate keys, each one established by some out-of-band procedure. A bank with a million customers needs a million keys, each handed over face to face. Scale that to a planet of strangers buying things from each other online, and the courier model collapses entirely. You cannot fly a briefcase to every website you have never visited before.

This chapter is about the idea that broke the problem open -- the idea that two people who have never met, communicating entirely over a wire that an eavesdropper reads in full, can nonetheless agree on a secret that the eavesdropper cannot compute. That sounds impossible: it is like two people shouting numbers across a crowded room and ending up with a secret nobody else in the room knows. Yet it can be done, it is not even hard, and the mathematics behind it is the same modular arithmetic we have been building since Chapter~3. The construction is Diffie--Hellman key exchange, and modern secure communication depends on it. When your browser negotiates a TLS connection, when two phones set up a Signal session, when an SSH client greets a server it has never spoken to, a descendant of Diffie--Hellman runs in the first few milliseconds.

We will build it the way we build everything: the intuition first, with a physical analogy and tiny numbers you can check by hand; then the precise protocol; then the hardness assumptions that make it secure and the careful distinction between two of them; then the attack that Diffie--Hellman does \emph{not} stop, which is the reason the rest of public-key cryptography exists; and finally a preview of forward secrecy, the property that makes the modern, ephemeral version of the protocol so valuable.

\chapterline

\section{The key-distribution problem}
\label{ch18:key-distribution}

As before, Alice wants to talk to Bob, and Eve sits on the wire and reads every bit that passes. The new wrinkle is that Alice and Bob have \emph{never met}. They share no key, no pad, no password. They have a public channel and nothing else.

With only the tools of the previous chapters, this is hopeless. Symmetric encryption is symmetric precisely because the same key locks and unlocks; if Alice has to send that key to Bob, she has to send it over the wire, and Eve reads it. Encrypting the key under another key only pushes the problem back a step: how did they agree on \emph{that} key? The recursion never bottoms out. Information-theoretically, if every bit Alice sends is visible to Eve, then anything Bob can compute from those bits, Eve can compute too. Bob has no advantage over Eve. They are looking at the same transcript.

\begin{intuitionbox}
The key-distribution problem in one sentence: Bob and Eve see exactly the same messages, so how can Bob end up knowing something Eve does not?

The resolution, which will take the rest of the section to make precise, is that Bob has one thing Eve lacks -- not a secret message on the wire, but a secret he \emph{chose} and never sent. Both Alice and Bob keep a private number to themselves and broadcast only a scrambled version of it. The key property is that their two private numbers can be combined into a shared value, but only by someone holding at least one of them. Eve holds neither.
\end{intuitionbox}

Before the 1970s, the standard count made the problem look unwinnable. A network of $n$ parties who all want to talk pairwise needs a distinct key for each pair, which is $\binom{n}{2} = n(n-1)/2$ keys. For $n = 1000$ that is just under half a million keys, every one of them established by some trusted, out-of-band method and stored securely forever. Add one new member and you must distribute $n$ fresh keys to the existing crowd. Militaries and embassies lived with this -- they had couriers and diplomatic pouches -- but it does not scale to commerce, and it certainly does not scale to the open internet.

\begin{historybox}
The breakthrough has a tangled history. In 1974 Ralph Merkle conceived ``puzzles'' that gave a quadratic advantage over an eavesdropper, and in 1976 Whitfield Diffie and Martin Hellman published ``New Directions in Cryptography,'' which laid out the public-key concept and gave the key-exchange protocol that now bears their names (Hellman has long insisted the method should be called Diffie--Hellman--Merkle). The paper opens with a line that reads, in hindsight, like a manifesto: ``We stand today on the brink of a revolution in cryptography.'' They were right, though open publication sparked a decade-long fight with the NSA over academic publication and export controls.

What no one outside the British signals intelligence agency GCHQ knew was that James Ellis, Clifford Cocks, and Malcolm Williamson had discovered ``non-secret encryption'' and a Diffie--Hellman-like exchange between 1969 and 1973. Their work was classified and stayed secret until 1997. So the idea was invented twice: once in the open, where it changed the world, and once in secret, where it sat in a vault. The open invention is the one that matters historically, because cryptography is only useful when people can actually use it.
\end{historybox}

\section{The intuition: mixing paint}
\label{ch18:paint}

The cleanest analogy for Diffie--Hellman uses paint, and it is worth internalizing before any algebra, because the algebra is just this picture with numbers.

Suppose mixing two colors of paint is easy but \emph{un}-mixing -- recovering the original colors from a blend -- is hard. Diffie--Hellman is built on that asymmetry.

\begin{enumerate}
  \item Alice and Bob publicly agree on a common starting color, say a particular yellow. Everyone, including Eve, sees this yellow.
  \item Alice picks a secret color, say red, and keeps it private. She mixes yellow~$+$~red into an orange blend and sends the orange to Bob over the public channel. Eve sees the orange but cannot un-mix it to recover Alice's red.
  \item Bob picks a secret color, say blue, keeps it private, mixes yellow~$+$~blue into a green blend, and sends the green to Alice. Eve sees the green but cannot extract Bob's blue.
  \item Alice stirs her secret red into the received green, while Bob stirs his secret blue into the received orange. Both arrive at the same three-color blend: \emph{yellow~$+$~red~$+$~blue}.
\end{enumerate}

Alice and Bob now share that final blend. Eve, watching the whole exchange, holds the public yellow, the orange, and the green. To reach the shared blend she would need to add either Alice's red or Bob's blue to something -- but those secret colors were never transmitted, only ever mixed in. Un-mixing is hard by assumption, so Eve is stuck with two mixtures she cannot pull apart.

\begin{figure}[H]
\centering
\begin{tikzpicture}[
  font=\small, >=Latex,
  box/.style={draw, rounded corners=2pt, minimum height=0.75cm, align=center},
  pub/.style={box, fill=quotientorange!20, minimum width=3.4cm},
  privA/.style={box, fill=aliceblue!20, minimum width=2.8cm},
  privB/.style={box, fill=bobgreen!20, minimum width=2.8cm},
  blendA/.style={box, fill=quotientorange!30, minimum width=2.9cm},
  blendB/.style={box, fill=targetgreen!25, minimum width=2.9cm},
  shared/.style={box, fill=structurepurple!18, font=\bfseries, minimum width=5.6cm}
]
  % Stage 1: Public Base
  \node[font=\footnotesize\scshape, text=gray] at (-2.9,3.4) {1. Common Base:};
  \node[pub] (base) at (2.8,3.4) {Public Common Base (Yellow $g$)};

  % Stage 2: Private Secrets
  \node[font=\footnotesize\scshape, text=gray] at (-2.9,2.1) {2. Secret Additions:};
  \node[privA] (a) at (0.3,2.1) {Alice: Secret Red ($a$)};
  \node[privB] (b) at (5.3,2.1) {Bob: Secret Blue ($b$)};

  % Stage 3: Public Mixtures
  \node[font=\footnotesize\scshape, text=gray] at (-2.9,0.6) {3. Transmitted Blends:};
  \node[blendA] (Ablend) at (0.3,0.6) {Orange ($A = g^a$)\\ \scriptsize (Yellow $+$ Red)};
  \node[blendB] (Bblend) at (5.3,0.6) {Green ($B = g^b$)\\ \scriptsize (Yellow $+$ Blue)};

  % Feed Yellow Base down the center gap to both blends
  \draw[->, thick, quotientorange!80!black] (base.south) -- (2.8,1.4) -| ([xshift=0.85cm]Ablend.north);
  \draw[->, thick, quotientorange!80!black] (base.south) -- (2.8,1.4) -| ([xshift=-0.85cm]Bblend.north);
  \draw[->, thick] (a) -- (Ablend);
  \draw[->, thick] (b) -- (Bblend);

  % Channel Exchange
  \draw[->, thick, trapred!80!black] (Ablend.south) -- ++(0,-0.45) -| (4.0,-0.55)
    node[pos=0.25, below, font=\scriptsize, text=gray] {public channel};
  \draw[->, thick, trapred!80!black] (Bblend.south) -- ++(0,-0.45) -| (1.6,-0.55);

  % Stage 4: Common Secret
  \node[font=\footnotesize\scshape, text=gray] at (-2.9,-1.35) {4. Shared Secret:};
  \node[shared] (s) at (2.8,-1.35) {Shared Secret: Yellow $+$ Red $+$ Blue ($s = g^{ab}$)};

  \draw[->, thick] (1.6,-0.55) -- (1.6,-0.98);
  \draw[->, thick] (4.0,-0.55) -- (4.0,-0.98);

  % Eve observer note
  \node[draw, rounded corners=2pt, dashed, trapred!80, fill=trapred!6, font=\scriptsize, align=center] at (8.3,0.6) {
    \textbf{Eve observes:}\\Yellow, Orange, Green.\\
    Cannot un-mix to get\\Red or Blue!
  };
\end{tikzpicture}
\caption{Diffie--Hellman represented as paint mixing. Both parties begin with a common public color (yellow $g$). Each party mixes in a private secret color and transmits the resulting mixture over the open wire. Adding their own private color to the received blend brings both parties to the identical shared three-color mixture ($g^{ab}$), while an eavesdropper cannot un-mix either public blend.}
\label{ch18:fig-paint}
\end{figure}

The remaining task is to find a numerical operation that behaves like paint: easy to mix, hard to un-mix, and one where the order of mixing does not matter, so that Alice's ``green then red'' equals Bob's ``orange then blue.'' Modular exponentiation is exactly that operation, and the order-independence is just the law of exponents.

\section{Diffie--Hellman key exchange}
\label{ch18:dh-protocol}

We work in a cyclic group, which for now you can read as ``the nonzero residues mod a prime, multiplied.'' Fix a large prime $p$ and a generator $g$ of the multiplicative group $\Zns{p}$. Recall from Chapter~5 that a generator is an element whose powers $g^1, g^2, g^3, \dots$ cycle through every nonzero residue mod $p$ before repeating; the group has order $p-1$, so $g^{p-1} \equiv 1 \pmod p$ and no smaller positive exponent does that. Both $p$ and $g$ are \emph{public}. Eve knows them. They can even be baked into a standard and shared by millions of connections.

``Mixing the public color with a secret'' becomes ``raising $g$ to a secret exponent.'' Here is the protocol.

\begin{focusbox}
\textbf{Diffie--Hellman key exchange.} Public parameters: a prime $p$ and a generator $g$ of $\Zns{p}$.
\begin{enumerate}
  \item Alice samples a secret exponent $a \getsr \{1, \dots, p-2\}$, computes $A = g^a \bmod p$, and sends $A$ to Bob.
  \item Bob samples a secret exponent $b \getsr \{1, \dots, p-2\}$, computes $B = g^b \bmod p$, and sends $B$ to Alice.
  \item Alice computes $s = B^a \bmod p$. Bob computes $s = A^b \bmod p$.
\end{enumerate}
Both arrive at the same value $s = g^{ab} \bmod p$, the shared secret.
\end{focusbox}

The correctness is one line of exponent algebra:
\[
  B^a = (g^b)^a = g^{ba} = g^{ab} = (g^a)^b = A^b \pmod p .
\]
Alice's secret $a$ plays the role of her red paint, Bob's secret $b$ is his blue, $g$ is the public yellow, and $g^a$, $g^b$ are the one-color blends sent across. The shared $g^{ab}$ is the three-color result, and the commutativity $g^{ba} = g^{ab}$ is precisely the order-independence the paint story needed. Eve sees $p$, $g$, $A = g^a$, and $B = g^b$. To finish, she would need $g^{ab}$, and the entire security of the scheme is the claim that recovering $g^{ab}$ from $g^a$ and $g^b$ is hard. We will name that claim shortly.

\begin{whybox}
Why can Eve not just read off $a$ from $A = g^a \bmod p$? Because inverting modular exponentiation -- finding the exponent $a$ given $g$, $p$, and $g^a \bmod p$ -- is the \emph{discrete logarithm problem}, and for well-chosen $p$ it is believed to be computationally infeasible. Computing $g^a \bmod p$ takes a few thousand multiplications via square-and-multiply (Chapter~3 and the companion algorithms book). Going the other way, with the best known general algorithms, takes time that grows sub-exponentially in the size of $p$ -- astronomically slow once $p$ has a couple thousand bits. This is the same one-way asymmetry as ``mixing is easy, un-mixing is hard.''
\end{whybox}

\subsection{A worked example with tiny numbers}
\label{ch18:tiny}

Real Diffie--Hellman uses primes thousands of bits long. To follow each step, we shrink everything to numbers you can verify on paper.

Take $p = 23$ and $g = 5$. First, a sanity check that $5$ generates $\Zns{23}$. The group has order $p - 1 = 22$. The order of $5$ must divide $22$, so it is one of $1, 2, 11, 22$. Squaring up: $5^2 = 25 \equiv 2$, $5^4 \equiv 4$, $5^8 \equiv 16 \pmod{23}$. Then $5^{11} = 5^8 \cdot 5^2 \cdot 5^1 \equiv 16 \cdot 2 \cdot 5 = 160 \equiv 160 - 6\cdot 23 = 22 \equiv -1 \pmod{23}$. Since $5^{11} \equiv -1 \neq 1$ and $5^2 \equiv 2 \neq 1$, the order is not $1, 2,$ or $11$; it must be $22$. So $5$ is a generator. Good.

Now run the protocol.

\begin{itemize}
  \item Alice picks $a = 6$. She computes $A = 5^6 \bmod 23$. Using $5^4 \equiv 4$ and $5^2 \equiv 2$, we get $5^6 = 5^4 \cdot 5^2 \equiv 4 \cdot 2 = 8$. Alice sends $A = 8$.
  \item Bob picks $b = 15$. He computes $B = 5^{15} \bmod 23$. Using $5^8 \equiv 16$, $5^4 \equiv 4$, $5^2 \equiv 2$, $5^1 \equiv 5$: $5^{15} = 5^8 \cdot 5^4 \cdot 5^2 \cdot 5^1 \equiv 16 \cdot 4 \cdot 2 \cdot 5$. Reduce as you go: $16 \cdot 4 = 64 \equiv 64 - 2\cdot 23 = 18$; $18 \cdot 2 = 36 \equiv 13$; $13 \cdot 5 = 65 \equiv 65 - 2\cdot 23 = 19$. Bob sends $B = 19$.
  \item Alice computes $s = B^a = 19^6 \bmod 23$. Note $19 \equiv -4$, so $19^6 \equiv (-4)^6 = 4^6$. Now $4^3 = 64 \equiv 18$, so $4^6 = (4^3)^2 = 18^2 = 324 \equiv 324 - 14\cdot 23 = 2$. So $s = 2$.
  \item Bob computes $s = A^b = 8^{15} \bmod 23$. Since $8 = 5^6$ we expect $g^{ab} = 5^{90}$, and $90 \bmod 22 = 90 - 4\cdot 22 = 2$, so $s = 5^2 = 2$. Direct check: $8^2 = 64 \equiv 18$, $8^4 = 18^2 = 324 \equiv 2$, $8^8 = 2^2 = 4$, and $8^{15} = 8^8 \cdot 8^4 \cdot 8^2 \cdot 8^1 \equiv 4 \cdot 2 \cdot 18 \cdot 8$. Reduce: $4\cdot 2 = 8$; $8 \cdot 18 = 144 \equiv 144 - 6\cdot 23 = 6$; $6 \cdot 8 = 48 \equiv 2$. So $s = 2$.
\end{itemize}

Both parties hold $s = 2 = g^{ab} \bmod 23$, computed by two different routes. Eve, watching the wire, knows $p = 23$, $g = 5$, $A = 8$, $B = 19$. To find $s$ she must solve a discrete log -- find $a$ with $5^a \equiv 8$, or $b$ with $5^b \equiv 19$. With $p = 23$ she can just try all $22$ exponents, and indeed she would find $a = 6$ in a moment. That is the point of a toy: the security vanishes at small sizes. Replace $23$ with a $2048$-bit prime and the exhaustive search becomes a search over $2^{2047}$ candidates, which no machine that will ever exist can finish.

\begin{figure}[H]
\centering
\begin{tikzpicture}[
  font=\small, >=Latex,
  box/.style={draw, rounded corners=2pt, minimum height=0.75cm, align=center},
  partyA/.style={box, fill=aliceblue!15, minimum width=3.4cm},
  partyB/.style={box, fill=bobgreen!15, minimum width=3.4cm},
  evebox/.style={draw, rounded corners=2pt, dashed, fill=evered!8, text=evered!90!black, font=\scriptsize, align=center, inner sep=5pt}
]
  % Party Headers
  \node[partyA] (ahdr) at (0,3.2) {\textbf{Alice}\\$p=23,\ g=5$};
  \node[partyB] (bhdr) at (7.2,3.2) {\textbf{Bob}\\$p=23,\ g=5$};

  % Alice local step 1
  \node[box, fill=aliceblue!8, minimum width=3.4cm] (astep1) at (0,1.9) {Secret $a = 6$\\$A = 5^6 \equiv 8 \pmod{23}$};

  % Bob local step 1
  \node[box, fill=bobgreen!8, minimum width=3.4cm] (bstep1) at (7.2,1.9) {Secret $b = 15$\\$B = 5^{15} \equiv 19 \pmod{23}$};

  % Shared secret computation
  \node[partyA, font=\bfseries] (astep2) at (0,-1.0) {Shared Key $s$:\\$s = B^a = 19^6 \equiv 2 \pmod{23}$};
  \node[partyB, font=\bfseries] (bstep2) at (7.2,-1.0) {Shared Key $s$:\\$s = A^b = 8^{15} \equiv 2 \pmod{23}$};

  % Local secret feed
  \draw[->, thick, aliceblue!80!black] (astep1.south) -- (astep2.north)
    node[midway, left=2pt, font=\scriptsize, text=aliceblue!80!black] {use $a=6$};
  \draw[->, thick, bobgreen!80!black] (bstep1.south) -- (bstep2.north)
    node[midway, right=2pt, font=\scriptsize, text=bobgreen!80!black] {use $b=15$};

  % Transmission arrows across wire
  \draw[->, thick, trapred!80!black] ([yshift=-0.1cm]astep1.east) -- ([yshift=0.15cm]bstep2.west)
    node[pos=0.35, above=3pt, font=\scriptsize, fill=white, inner sep=1.5pt, text=trapred!85!black] {public wire: $A = 8$};
  \draw[->, thick, trapred!80!black] ([yshift=-0.1cm]bstep1.west) -- ([yshift=0.15cm]astep2.east)
    node[pos=0.65, below=3pt, font=\scriptsize, fill=white, inner sep=1.5pt, text=trapred!85!black] {public wire: $B = 19$};

  % Eve observer box
  \node[evebox] at (3.6,-2.3) {
    \textbf{Eve Observes Public Transcript:}\\
    $p = 23,\ g = 5,\ A = 8,\ B = 19$\\
    $\implies$ To compute $s = 2$, Eve must solve $5^a \equiv 8 \pmod{23}$ (Discrete Logarithm).
  };
\end{tikzpicture}
\caption{The concrete Diffie--Hellman exchange with toy numbers $p=23, g=5$. Alice and Bob exchange public powers $A = 8$ and $B = 19$ across the insecure channel. Each computes the shared secret $s = 2$ using their private exponent. Eve must solve the discrete logarithm problem to reconstruct $s$.}
\label{ch18:fig-flow}
\end{figure}

\begin{implementbox}
A toy Diffie--Hellman in Python. This is for \textbf{learning only} -- it uses a microscopic prime, samples exponents with a non-cryptographic RNG, and derives a key by naively hashing the group element. Never use anything like this for real traffic; reach for a vetted library (\texttt{cryptography}, \texttt{libsodium}) that uses standardized parameters and constant-time arithmetic.
\begin{lstlisting}[language=Python]
import random, hashlib

p = 23          # toy prime -- a real one is >= 2048 bits
g = 5           # generator of (Z/23Z)*

def keypair():
    x = random.randrange(2, p - 1)   # secret exponent (use os.urandom in reality!)
    X = pow(g, x, p)                 # public value g^x mod p, via fast modexp
    return x, X

a, A = keypair()                     # Alice
b, B = keypair()                     # Bob

# each side raises the OTHER's public value to its OWN secret
s_alice = pow(B, a, p)
s_bob   = pow(A, b, p)
assert s_alice == s_bob              # both get g^(ab) mod p

# turn the shared group element into a symmetric key
shared = s_alice
key = hashlib.sha256(str(shared).encode()).digest()
print("shared group element:", shared)
print("derived key (hex):", key.hex())
\end{lstlisting}
Note the symmetry of \texttt{pow(B, a, p)} and \texttt{pow(A, b, p)}: each party exponentiates the value it received by the secret it kept. Apart from choosing the parameters, that is the entire protocol.
\end{implementbox}

\begin{trapbox}
Three pitfalls already visible in the toy code, all fatal in practice.

\textbf{Tiny or weak parameters.} A small $p$, or a $p$ where $p-1$ has only small prime factors, lets attackers solve the discrete log fast (the Pohlig--Hellman attack, Chapter~20, decomposes the problem across the small factors of $p-1$). Real deployments use large \emph{safe primes} where $p = 2q+1$ for a prime $q$.

\textbf{Bad randomness.} \texttt{random.randrange} is a Mersenne Twister, not cryptographic; an attacker who can predict it recovers your secret exponent and the whole exchange falls. Use the operating system CSPRNG (Chapter~17).

\textbf{Using the raw group element as a key.} $g^{ab} \bmod p$ is an integer in a structured set, not a uniform bitstring. You must run it through a key-derivation function (a hash or HKDF) before using it as an AES or MAC key. The toy hashes it; production code uses HKDF with a salt and context label.
\end{trapbox}

\section{What ``secure'' means here: three hardness assumptions}
\label{ch18:assumptions}

So far ``hard to un-mix'' has been an informal analogy. To build a security proof on top of it, we need to state that hardness assumption precisely. Three related problems underlie Diffie--Hellman, and the differences between them matter. We state each in a cyclic group $\G = \langle g \rangle$ of order $q$ -- for the protocol above, $\G = \Zns{p}$ and $q = p-1$, though we will see later that prime-order subgroups and elliptic-curve groups are better choices.

\begin{definition}[Discrete logarithm problem, DLP]
\label{ch18:dlp}
Given $g$ and $h = g^x$ in $\G$, find $x \in \{0, 1, \dots, q-1\}$. The DLP is \emph{hard} in $\G$ if every $\PPT$ algorithm solves it with only negligible probability.
\end{definition}

\begin{definition}[Computational Diffie--Hellman, CDH]
\label{ch18:cdh}
Given $g$, $g^a$, and $g^b$ in $\G$ (with $a, b$ uniform and unknown), compute $g^{ab}$. The CDH assumption says every $\PPT$ algorithm succeeds with only negligible probability.
\end{definition}

\begin{definition}[Decisional Diffie--Hellman, DDH]
\label{ch18:ddh}
Given $g$, $g^a$, $g^b$, and a fourth element $z$ that is \emph{either} $g^{ab}$ \emph{or} a uniformly random group element $g^c$ (each with probability $1/2$), decide which. The DDH assumption says no $\PPT$ algorithm distinguishes the two cases with advantage better than negligible:
\[
  \abs[\big]{\Pr[\,\mathcal{A}(g, g^a, g^b, g^{ab}) = 1\,] - \Pr[\,\mathcal{A}(g, g^a, g^b, g^c) = 1\,]} \le \negl .
\]
\end{definition}

These line up on a ladder of strength. If you can solve DLP you can solve CDH: recover $a$ from $g^a$, then raise $g^b$ to the $a$ to get $g^{ab}$. And if you can solve CDH you can solve DDH: compute $g^{ab}$ yourself and check whether $z$ equals it. So
\[
  \text{DLP hard} \;\Longleftarrow\; \text{CDH hard} \;\Longleftarrow\; \text{DDH hard},
\]
reading ``$X \Leftarrow Y$'' as ``$Y$ being hard is at least as strong an assumption as $X$ being hard.'' DDH is the strongest assumption (easiest to break), DLP the weakest (hardest to break). Whether the implications reverse -- whether breaking DLP is the \emph{only} way to break CDH, for instance -- is not known. In some groups they might come apart.

\begin{whybox}
Why do we need DDH and not just CDH? Because the goal of key exchange is not merely that Eve cannot \emph{compute} the whole shared secret, but that the shared secret \emph{looks random} to her -- so that we can safely feed it into a hash and get a key with no exploitable structure. CDH only promises Eve cannot produce $g^{ab}$ outright. DDH promises something stronger: $g^{ab}$ is indistinguishable from a random group element even to an Eve who only wants to learn one bit about it. That indistinguishability is what lets us treat the derived key as good as random in a security proof. In short, CDH says Eve cannot compute the secret; DDH says she cannot even distinguish it from random, and security proofs need the latter.
\end{whybox}

\begin{trapbox}
DDH is genuinely \emph{false} in the full group $\Zns{p}$, and the reason is a short attack. The map $x \mapsto x^{(p-1)/2} \bmod p$ -- the Legendre symbol -- reveals whether $x$ is a quadratic residue (a perfect square mod $p$), outputting $+1$ or $-1$. The QR-ness of $g^{ab}$ is a deterministic function of the QR-ness of $g^a$ and $g^b$, which Eve can read off directly. So Eve computes the residue character that the true $g^{ab}$ must have and checks whether $z$ matches it; a random $g^c$ fails the check half the time. That single leaked bit breaks DDH.

The fix is to work in a \emph{prime-order subgroup}: pick a safe prime $p = 2q+1$ and run the protocol in the order-$q$ subgroup of quadratic residues, where every element is a QR and the Legendre attack says nothing. This is exactly why standardized finite-field DH parameters use safe primes and restrict to the prime-order subgroup, and why elliptic-curve groups (Chapter~21), which have prime order by construction, are now preferred -- DDH is believed to hold there cleanly.
\end{trapbox}

\begin{connectionbox}
The cost asymmetry that makes all of this work is quantified in the companion algorithms book. Computing $g^a \bmod p$ by square-and-multiply costs $O(\log a)$ modular multiplications -- a few thousand for a $2048$-bit exponent. The best known classical attack on finite-field discrete log, the general number field sieve, runs in time roughly $\exp\!\big(c\,(\ln p)^{1/3}(\ln \ln p)^{2/3}\big)$: sub-exponential, but still wildly out of reach for a $2048$-bit prime. On elliptic curves even that sieve does not apply, and the best attacks are the generic square-root methods (Pollard's rho), costing about $\sqrt{q}$ group operations -- which is why a $256$-bit elliptic curve buys security comparable to a $3072$-bit finite-field prime. The probability and number-theory chapters of this book supply the group theory; the algorithms book supplies the running-time bounds.
\end{connectionbox}

\subsection{From shared secret to usable key}
\label{ch18:kdf}

The shared value $s = g^{ab}$ is an element of $\G$, not a key. Three things have to happen before it can encrypt anything. First, it must be encoded as a fixed-length byte string. Second, even granting DDH, $s$ is uniform over $\G$, not over all bitstrings of its length -- a $2048$-bit prime leaves gaps in the $2048$-bit space. Third, you usually want several keys from one exchange (one for each direction, perhaps separate keys for encryption and MAC). The standard remedy for all three is a key-derivation function, almost always HKDF (Chapter~15), used as
\[
  k_1 \concat k_2 \concat \cdots = \mathrm{HKDF}\big(s,\ \text{salt},\ \text{context}\big),
\]
where the context label binds the keys to this protocol and these parties. Feeding $s$ straight into AES is a real-world bug; the KDF step is required.

\section{The man-in-the-middle problem}
\label{ch18:mitm}

Plain Diffie--Hellman has a serious weakness, and the certificates and signatures of later chapters exist largely to deal with it. Plain Diffie--Hellman defeats a \emph{passive} eavesdropper who only reads the wire. It does nothing against an \emph{active} attacker who can intercept and replace messages. The reason is simple: nothing in the protocol ties the value $A$ to Alice or $B$ to Bob. They are just numbers. Eve can substitute her own.

\begin{attackbox}
\textbf{Man-in-the-middle on unauthenticated Diffie--Hellman.}

Eve sits between Alice and Bob and relays their connection, but tampers. She runs two separate exchanges, pretending to be Bob toward Alice and Alice toward Bob.

\begin{enumerate}
  \item Alice sends $A = g^a$, intending it for Bob. Eve intercepts it. Bob never sees it.
  \item Eve picks her own secret $e$, computes $E = g^e$, and sends $E$ to Bob as if it came from Alice. She also sends $E$ to Alice as if it came from Bob.
  \item Bob, seeing $E$, computes his shared secret $s_B = E^b = g^{eb}$, believing he shares it with Alice. Bob sends $B = g^b$ back toward Alice; Eve intercepts that too.
  \item Alice, seeing $E$, computes $s_A = E^a = g^{ea}$, believing she shares it with Bob.
  \item Eve computes both: $s_A = A^e = g^{ae}$ and $s_B = B^e = g^{be}$.
\end{enumerate}

Now Eve shares one key with Alice ($g^{ae}$) and a different key with Bob ($g^{be}$). When Alice encrypts a message under $s_A$ and sends it ``to Bob,'' Eve decrypts it with $s_A$, reads it, re-encrypts it under $s_B$, and forwards it to Bob. Neither Alice nor Bob notices anything wrong: each completed a perfect Diffie--Hellman exchange and holds a working key. They are simply both talking to Eve. Confidentiality is gone, and so is integrity -- Eve can alter messages at will before re-encrypting.
\end{attackbox}

\begin{figure}[H]
\centering
\begin{tikzpicture}[
  font=\small, >=Latex,
  box/.style={draw, rounded corners=2pt, minimum height=0.72cm, align=center},
  partyA/.style={box, fill=aliceblue!15, minimum width=2.8cm},
  partyB/.style={box, fill=bobgreen!15, minimum width=2.8cm},
  mitm/.style={box, fill=trapred!18, draw=trapred!80!black, minimum width=3.5cm}
]
  % Headers
  \node[partyA] (ahdr) at (0,2.6) {\textbf{Alice}\\ \scriptsize Believes talking to Bob};
  \node[mitm] (ehdr) at (4.6,2.6) {\textbf{Eve (Active MITM)}\\ \scriptsize Intercepts \& Impersonates};
  \node[partyB] (bhdr) at (9.2,2.6) {\textbf{Bob}\\ \scriptsize Believes talking to Alice};

  % Alice to Eve
  \draw[->, thick, trapred!80!black] (0,1.5) -- node[above, font=\scriptsize] {$A = g^a$} (2.8,1.5);
  \node[font=\tiny, text=trapred!90!black] at (2.8, 1.8) {intercepted!};

  % Eve to Bob
  \draw[->, thick, trapred!80!black] (6.4,1.5) -- node[above, font=\scriptsize] {forged $E = g^e$} (9.2,1.5);

  % Bob to Eve
  \draw[->, thick, trapred!80!black] (9.2,0.5) -- node[above, font=\scriptsize] {$B = g^b$} (6.4,0.5);
  \node[font=\tiny, text=trapred!90!black] at (6.4, 0.8) {intercepted!};

  % Eve to Alice
  \draw[->, thick, trapred!80!black] (2.8,0.5) -- node[above, font=\scriptsize] {forged $E = g^e$} (0,0.5);

  % Key derivations
  \node[partyA, font=\footnotesize] at (0,-0.8) {Derived Key:\\$s_A = E^a = g^{ea}$};
  \node[mitm, font=\footnotesize] at (4.6,-0.8) {Eve computes BOTH keys:\\$s_A = A^e = g^{ae}$\\$s_B = B^e = g^{be}$};
  \node[partyB, font=\footnotesize] at (9.2,-0.8) {Derived Key:\\$s_B = E^b = g^{eb}$};

  % Traffic relay note
  \node[draw, rounded corners=2pt, dashed, trapred!80, fill=trapred!6, font=\scriptsize, align=center, inner sep=4pt] at (4.6,-2.2) {
    \textbf{Transparent Man-in-the-Middle Relay:}\\
    Alice sends $\Enc_{s_A}(m) \implies$ Eve decrypts with $s_A$, alters $m \to m'$, sends $\Enc_{s_B}(m') \implies$ Bob decrypts with $s_B$.\\
    Unauthenticated Diffie--Hellman provides no identity verification!
  };
\end{tikzpicture}
\caption{The man-in-the-middle attack on unauthenticated Diffie--Hellman. Eve intercepts public values $A$ and $B$, substituting her own value $E = g^e$. She completes independent exchanges with Alice and Bob, establishing distinct shared keys $s_A$ and $s_B$. Eve transparently decrypts, reads, alters, and re-encrypts all forwarded traffic.}
\label{ch18:fig-mitm}
\end{figure}

The lesson is that key exchange and \emph{authentication} are different problems, and Diffie--Hellman solves only the first. The exchange establishes a shared secret with \emph{someone}; it gives no assurance about \emph{who} that someone is. To pin down identity you need an external anchor: something Alice already knows about Bob (a public key she trusts) or a trusted third party who vouches for Bob's public value. That is the job of digital signatures (Chapter~22), certificates, and the public-key infrastructure (Chapter~24). In a real TLS handshake the server signs its Diffie--Hellman value with the private key matching the certificate your browser validates, so a substituted $E$ would carry an invalid signature and the connection would abort. Without that signature, Diffie--Hellman alone is wide open to the attack above.

\begin{intuitionbox}
A useful slogan: \emph{Diffie--Hellman gets you a secure pipe to somebody; authentication tells you who is on the other end.} Both are necessary. Encrypting to the wrong party with perfect secrecy is no comfort at all.
\end{intuitionbox}

\section{Ephemeral Diffie--Hellman and forward secrecy}
\label{ch18:forward-secrecy}

One more property explains why Diffie--Hellman -- rather than, say, encrypting a fresh key under the server's RSA public key -- is at the center of modern TLS. It concerns what happens \emph{after} a conversation, when the secrets get stolen.

Distinguish two ways to use the protocol. In \emph{static} Diffie--Hellman, a party reuses a long-term secret exponent across many sessions: the server's $b$ is fixed, baked into its key material, used for every client. In \emph{ephemeral} Diffie--Hellman (written DHE for finite fields, ECDHE for elliptic curves), each party generates a \emph{fresh} random exponent for every single session and throws it away the moment the session key is derived.

The payoff is \emph{forward secrecy} (also called perfect forward secrecy): the compromise of long-term secrets in the future does not expose past sessions. Picture an adversary who records all of Alice and Bob's encrypted traffic for years, patiently storing ciphertext she cannot yet read. Later she breaks into Bob's server and steals every long-term private key it holds. What can she decrypt retroactively?

With ephemeral Diffie--Hellman: nothing (Figure~\ref{ch18:fig-forward-secrecy}). The session key for a past conversation was $g^{ab}$ for that session's $a$ and $b$, both of which were deleted right after use. The long-term key she stole was used only to \emph{sign} the ephemeral values for authentication, never to compute the shared secret. The recorded ciphertext stays opaque forever.

With static key transport -- the old TLS mode where the client encrypts the session key under the server's long-term RSA public key -- the same break is catastrophic. The session key was protected solely by the server's long-term private key, so stealing that key later lets the adversary decrypt every recorded session that ever used it. Years of recorded traffic become readable at once.

\begin{figure}[H]
\centering
\begin{tikzpicture}[
  font=\small, >=Latex,
  box/.style={draw, rounded corners=2pt, minimum height=0.72cm, align=center, inner sep=4pt},
  sess/.style={box, minimum width=2.8cm, font=\footnotesize},
  breach/.style={box, fill=trapred!18, draw=trapred!80!black, minimum width=2.8cm, font=\footnotesize\bfseries}
]
  % Panel 1: Static Key Transport
  \node[font=\footnotesize\bfseries, text=trapred!85!black] at (-5.6, 2.5) {Static Key Transport:};
  \node[font=\scriptsize, text=gray] at (-5.6, 2.05) {(No Forward Secrecy)};

  \node[sess, fill=trapred!10, draw=trapred!70!black] (st1) at (-2.0, 2.3) {Session 1 ($t_1$)\\ \scriptsize $\Enc_{\mathsf{pk}_{\mathrm{server}}}(k_1)$};
  \node[sess, fill=trapred!10, draw=trapred!70!black] (st2) at (1.5, 2.3) {Session 2 ($t_2$)\\ \scriptsize $\Enc_{\mathsf{pk}_{\mathrm{server}}}(k_2)$};
  \node[breach] (stB) at (5.0, 2.3) {Server Breach ($t_3$)\\ \scriptsize Steals long-term $\mathsf{sk}$};

  \draw[->, thick, dashed, trapred!80!black] (stB.north) -- ++(0,0.45) -| (st1.north)
    node[pos=0.25, above=1pt, font=\scriptsize, text=trapred!85!black] {retroactive decryption unravels all past recorded traffic};
  \draw[->, thick, dashed, trapred!80!black] (stB.north) -- ++(0,0.45) -| (st2.north);

  % Divider
  \draw[dashed, gray!50] (-7.1, 1.05) -- (6.5, 1.05);

  % Panel 2: Ephemeral Diffie-Hellman (DHE / ECDHE)
  \node[font=\footnotesize\bfseries, text=targetgreen!80!black] at (-5.6, 0.0) {Ephemeral (EC)DHE:};
  \node[font=\scriptsize, text=gray] at (-5.6, -0.45) {(Forward Secrecy)};

  \node[sess, fill=targetgreen!12, draw=targetgreen!70!black] (ep1) at (-2.0, -0.2) {Session 1 ($t_1$)\\ \scriptsize $k_1 = g^{a_1 b_1}$ ($a_1,b_1$ erased)};
  \node[sess, fill=targetgreen!12, draw=targetgreen!70!black] (ep2) at (1.5, -0.2) {Session 2 ($t_2$)\\ \scriptsize $k_2 = g^{a_2 b_2}$ ($a_2,b_2$ erased)};
  \node[breach] (epB) at (5.0, -0.2) {Server Breach ($t_3$)\\ \scriptsize Steals signing $\mathsf{sk}_{\mathrm{sign}}$};

  \node[font=\scriptsize, text=targetgreen!80!black] at (1.5, -1.05) {$\mathsf{sk}_{\mathrm{sign}}$ only authenticated transcripts; erased $a_i, b_i$ cannot be recovered $\implies$ past sessions stay safe};
\end{tikzpicture}
\caption{Why ephemeral Diffie--Hellman provides forward secrecy. Under static key transport (top), stealing the server's long-term decryption key $\mathsf{sk}$ at time $t_3$ lets an eavesdropper retroactively decrypt every recorded session ($t_1, t_2$). Under ephemeral (EC)DHE (bottom), each session derives $k_i = g^{a_i b_i}$ from one-time exponents that are immediately erased; stealing the server's long-term signing key $\mathsf{sk}_{\mathrm{sign}}$ at $t_3$ reveals nothing about past session keys.}
\label{ch18:fig-forward-secrecy}
\end{figure}

\begin{whybox}
Forward secrecy works because of a deliberate separation of duties. The long-term key authenticates -- it answers ``who are you?'' -- and only that. The ephemeral exchange establishes confidentiality, using secrets that exist for one session and are then destroyed. Because the confidentiality secret is never derivable from any stored long-term value, there is nothing left for a future thief to use against an old recording. You cannot un-delete $a$ and $b$; they are gone. The security of last week's call no longer depends on a key that is still sitting on a server waiting to be stolen.
\end{whybox}

\begin{applicationbox}
Ephemeral Diffie--Hellman is everywhere in the protocols you use daily.

\textbf{TLS 1.3} (Chapter~25) removed static RSA key transport entirely. Every TLS~1.3 handshake uses ephemeral (EC)DHE, so forward secrecy is mandatory, not optional. The client offers a fresh ECDHE public value in its very first message; the server replies with its own and signs the transcript with its certificate key.

\textbf{Signal} (Chapter~27) goes further. Its X3DH handshake combines several Diffie--Hellman exchanges -- some ephemeral, some using medium-term ``prekeys'' -- and the Double Ratchet then performs a fresh ECDH on (almost) every message, continually replacing keys. This gives both forward secrecy and \emph{post-compromise security}: even if a device is briefly compromised, future messages heal back to secure once the attacker loses access.

\textbf{SSH} (Chapter~28) and \textbf{IPsec/IKE} (Chapter~30) both run ephemeral Diffie--Hellman to set up each session, authenticated by host keys or certificates.

The standardized parameters matter too. Finite-field DHE uses the safe primes of RFC~7919; elliptic-curve ECDHE overwhelmingly uses Curve25519, whose careful design makes the implementation pitfalls of this chapter (small subgroups, bad parameters, non-constant-time arithmetic) hard to trip over.
\end{applicationbox}

\begin{connectionbox}
The phrase ``perfect forward secrecy'' predates its TLS fame; it comes from 1990s work on authenticated key exchange (the STS protocol of Diffie, van Oorschot, and Wiener, and later the IKE design). For a long time forward secrecy was treated as a nice-to-have that cost extra computation, and many servers ran static RSA for speed. The 2013 Snowden disclosures, which revealed bulk recording of internet traffic, quickly changed this view: ``record now, decrypt later'' stopped being hypothetical. Within a few years ephemeral Diffie--Hellman went from optional to default to, in TLS~1.3, the only choice. The same ``harvest now, decrypt later'' worry now drives post-quantum migration (Chapter~35).
\end{connectionbox}

\section{Looking ahead}
\label{ch18:ahead}

Diffie--Hellman gave us the first piece of public-key cryptography: two strangers agreeing on a secret over an open line. But it left two threads dangling that the next chapters pick up. The man-in-the-middle attack showed that key exchange without authentication fails against an active attacker, which points to digital signatures and certificates. And the discrete-log assumption that secured the exchange turns out to be the same hard problem behind ElGamal encryption and the DSA/ECDSA signature schemes -- so the structure here recurs. The very next chapter takes a different hard problem, integer factoring, and builds RSA: the other foundational public-key system, which can encrypt and sign directly rather than only exchange keys. With both DLP-based and factoring-based cryptography in hand, the machinery of TLS, Signal, and SSH can be built from parts you understand.

\section{Exercises}
\label{ch18:exercises}

\begin{exercise}
With $p = 23$ and $g = 5$ as in the chapter, suppose Alice's secret is $a = 4$ and Bob's is $b = 3$. Compute $A = g^a \bmod 23$, $B = g^b \bmod 23$, and the shared secret $s = g^{ab} \bmod 23$ two ways (as $B^a$ and as $A^b$), checking that they agree.
\end{exercise}

\begin{exercise}
Take $p = 11$. Show that $g = 2$ is a generator of $\Zns{11}$ by listing the powers $2^1, 2^2, \dots, 2^{10} \bmod 11$ and confirming they hit all of $\{1, \dots, 10\}$. Then show that $g = 3$ is \emph{not} a generator by finding the order of $3$ mod $11$.
\end{exercise}

\begin{exercise}
(Breaking a toy.) Eve observes a Diffie--Hellman exchange with $p = 23$, $g = 5$, $A = 4$, and $B = 16$. By trying exponents (or by any method you like), recover Alice's secret $a$ and Bob's secret $b$, then compute the shared secret $s$. How many exponentiations did your search take in the worst case, and how does that count scale with the size of $p$?
\end{exercise}

\begin{exercise}
Prove the reduction CDH $\le$ DLP: given a $\PPT$ algorithm that solves the discrete logarithm problem in $\G$, construct one that solves CDH in $\G$ with the same success probability. Then prove DDH $\le$ CDH similarly. State clearly why neither reduction directly runs in the other direction.
\end{exercise}

\begin{exercise}
(The Legendre attack on DDH in $\Zns{p}$.) Let $p = 23$, so $g = 5$ is a generator and the quadratic residues mod $23$ are the even powers of $5$. Recall that $x$ is a QR iff $x^{11} \equiv 1 \pmod{23}$. Given a DDH challenge $(g^a, g^b, z)$, explain how to compute the QR character that a genuine $g^{ab}$ must have from the characters of $g^a$ and $g^b$, and how matching $z$ against it distinguishes a real $g^{ab}$ from a random group element with probability $3/4$ over a single challenge. Why does restricting to the subgroup of quadratic residues kill this attack?
\end{exercise}

\begin{exercise}
Carry out the man-in-the-middle attack of Section~\ref{ch18:mitm} concretely with $p = 23$, $g = 5$. Let Alice choose $a = 6$, Bob choose $b = 15$, and Eve choose $e = 9$. Compute the public value $E$ Eve sends to both, the key $s_A = g^{ae}$ Eve shares with Alice, and the key $s_B = g^{be}$ Eve shares with Bob. Verify that Alice and Bob each derive a key that matches one of Eve's.
\end{exercise}

\begin{exercise}
Explain in one or two paragraphs why hashing the shared group element through a KDF is necessary even when DDH holds. Give a specific example of something that could go wrong if you used the raw integer $g^{ab} \bmod p$ directly as an AES-128 key (consider its length and its distribution).
\end{exercise}

\begin{exercise}
A server runs \emph{static} Diffie--Hellman: it uses the same secret exponent $b$ for every client, publishing $B = g^b$ as a long-lived public value. An attacker records a year of encrypted sessions and later steals $b$. Show that the attacker can now recompute the shared secret of every recorded session. Then explain precisely which step of the ephemeral variant prevents this, and why deleting the ephemeral exponents is the essential step.
\end{exercise}

\begin{exercise}
(Small-subgroup pitfall, preview of Chapter~21.) Suppose $p = 23$ so $\Zns{p}$ has order $22 = 2 \cdot 11$. The element $h = 22 \equiv -1$ has order $2$. Imagine Bob fails to validate Alice's incoming value and Eve sends $A = 22$ instead of a legitimate public value. Show that Bob's computed shared secret $A^b \bmod 23$ can then take only one of two values depending on the parity of $b$, leaking a bit of $b$. What general validation rule on incoming group elements prevents this class of attack?
\end{exercise}

\begin{exercise}
Estimate, using the cost figures from Section~\ref{ch18:assumptions}, how the work of the legitimate parties compares to the work of an attacker for a $2048$-bit prime. Roughly how many modular multiplications does Alice perform to compute $g^a \bmod p$? Contrast that with the (sub-exponential) cost of the number field sieve, and explain in words why this gap is what makes the scheme usable.
\end{exercise}
